\subsection{Unión-Find (Conjuntos Disjuntos) para Componentes Conexas}

\graybold{Preguntas guía}

\begin{itemize}
\item ¿Necesito agrupar elementos en conjuntos según relaciones que se van descubriendo una a una (por ejemplo, aristas de un grafo)?
\item ¿El grafo es no dirigido y solo me interesa saber qué nodos quedan conectados entre sí, sin importar el camino ni el peso?
\item ¿Voy a intercalar muchas operaciones de "unir dos elementos" con consultas de "¿están en el mismo grupo?"?
\item ¿Necesito contar cuántas componentes conexas resultan al final, incluyendo nodos que quedan aislados?
\end{itemize}

\graybold{¿Para qué sirve?} La estructura Unión-Find (o DSU, Disjoint Set Union) mantiene una partición de un conjunto de elementos en subconjuntos disjuntos, permitiendo unir dos subconjuntos y consultar si dos elementos pertenecen al mismo subconjunto en tiempo casi constante amortizado. Es la herramienta estándar para detectar y contar componentes conexas de un grafo no dirigido sin necesidad de recorrerlo explícitamente con BFS/DFS, y en general para cualquier problema de agrupamiento incremental por relaciones de equivalencia.

\graybold{Complejidad:} \bigO{n \alpha(n)} amortizado para $n$ operaciones de unión/búsqueda, donde $\alpha$ es la inversa de la función de Ackermann (en la práctica, una constante menor a 5 para cualquier entrada realista).

\graybold{Fundamento teórico:} Cada elemento apunta a un "padre" dentro de un bosque de árboles; la raíz de cada árbol identifica a la componente a la que pertenece. La operación \texttt{find} sigue punteros de padre hasta la raíz, y se le aplica \emph{compresión de camino}: en cada búsqueda, todos los nodos visitados se reenganchan directamente a la raíz, aplanando el árbol para futuras consultas. La operación \texttt{union} usa además \emph{unión por tamaño}: la raíz del árbol más pequeño se cuelga bajo la raíz del árbol más grande, evitando que los árboles crezcan en profundidad de forma innecesaria. La combinación de ambas heurísticas es lo que garantiza la cota $\bigO{n \alpha(n)}$; usando solo una de las dos, la cota se degrada a $\bigO{n \log n}$. Al final, el número de raíces distintas entre todos los nodos del universo (aquí, las letras de \texttt{A} hasta la letra máxima dada) es exactamente el número de componentes conexas maximales, ya que cada nodo sin aristas es trivialmente su propia componente de tamaño 1.

\graybold{Ejercicio aplicado}

Dado un grafo no dirigido cuyos nodos son letras mayúsculas de \texttt{A} hasta una letra máxima dada (todas existen, tengan o no aristas), y una lista de aristas, calcular cuántos subgrafos conexos maximales (componentes conexas) tiene.

\graybold{I/O:} Se usó lectura línea a línea (\texttt{sys.stdin.buffer.readline}) en vez de tokenización total, porque la estructura del problema depende de las líneas en blanco: estas delimitan tanto el bloque inicial (cantidad de casos) como el final de la lista de aristas de cada caso, información que se perdería al tokenizar todo con \texttt{.split()}. Los límites son mínimos (a lo sumo 26 nodos por caso, ya que son letras \texttt{A--Z}), por lo que no se requirió ninguna optimización de rendimiento; el foco estuvo en la robustez del parseo (múltiples líneas en blanco consecutivas, ausencia de línea en blanco final antes de EOF), verificada con varios casos de prueba construidos a mano además del ejemplo dado.

\begin{lstlisting}[language=Python]
import sys

def main():
    read = sys.stdin.buffer.readline

    def next_line():
        # Lee una linea cruda y la decodifica; None indica EOF
        line = read()
        if not line:
            return None
        return line.decode().rstrip('\n\r')

    # Saltar posibles lineas en blanco antes del contador de casos
    line = next_line()
    while line is not None and line.strip() == '':
        line = next_line()
    if line is None:
        return
    n = int(line.strip())

    resultados = []
    for _ in range(n):
        # Saltar lineas en blanco entre casos (puede haber mas de una)
        line = next_line()
        while line is not None and line.strip() == '':
            line = next_line()

        max_letra = line.strip()
        # Todos los nodos de A a max_letra existen, tengan o no aristas
        padre = {chr(c): chr(c) for c in range(ord('A'), ord(max_letra) + 1)}
        tam = {c: 1 for c in padre}

        def encontrar(x):
            raiz = x
            while padre[raiz] != raiz:
                raiz = padre[raiz]
            # Compresion de camino: reengancha todo el trayecto a la raiz
            while padre[x] != raiz:
                padre[x], x = raiz, padre[x]
            return raiz

        def unir(a, b):
            ra, rb = encontrar(a), encontrar(b)
            if ra == rb:
                return
            # Union por tamano: el arbol chico cuelga del grande
            if tam[ra] < tam[rb]:
                ra, rb = rb, ra
            padre[rb] = ra
            tam[ra] += tam[rb]

        # Leer aristas hasta encontrar linea en blanco o EOF
        while True:
            line = next_line()
            if line is None or line.strip() == '':
                break
            arista = line.strip()
            a, b = arista[0], arista[1]
            unir(a, b)

        raices = set(encontrar(nodo) for nodo in padre)
        resultados.append(len(raices))

    sys.stdout.write('\n\n'.join(str(r) for r in resultados) + '\n')

main()
\end{lstlisting}